\documentclass[10pt]{article}
\usepackage[letterpaper,margin=0.45in]{geometry}
\usepackage{booktabs}
\usepackage{graphicx}
\usepackage{amsmath}
\makeatletter\setlength{\@fptop}{0pt}\makeatother
\begin{document}
\begin{table}[!t]
\centering
\caption{Deployment efficiency for serving four requested horizons. Values are macro-averaged over seven datasets on one exclusive RTX 3090 with FP32 and batch size 1.}
\label{tab:iscf-bsca-efficiency}
\scriptsize
\resizebox{\textwidth}{!}{%
\begin{tabular}{lrrrrrrrl}
\toprule
System & Models & Params (M) & Ckpt. (MiB) & Train (GPU h) & Single (ms) & All-$H$ (ms) & Peak (MiB) & CHPC \\
\midrule
\textbf{ISCF-BSCA} & 1 & 2.926 & 17.68 & 2.028 & 10.306 & 10.318 & 38.8 & Architectural \\
TimeAlign & 4 & 10.741 & 95.43 & 0.290 & 0.928 & 3.582 & 109.1 & No \\
QDF & 4 & 5.337 & 20.38 & 0.204 & 1.141 & 4.426 & 31.9 & No \\
DLinear-$H720$-prefix & 1 & 0.485 & 1.85 & 0.021 & 0.405 & 0.429 & 11.1 & Service protocol \\
PatchTST-$H720$-prefix & 1 & 3.198 & 12.23 & 1.110 & 3.933 & 3.882 & 58.0 & Service protocol \\
\bottomrule
\end{tabular}%
}
\vspace{2pt}
\parbox{\textwidth}{\footnotesize Models, parameters, checkpoints and training time count the complete deployed service per dataset. Single is the arithmetic mean over requests $H\in\{96,192,336,720\}$. All-$H$ uses one $H=720$ forward plus prefix views for one-model services and four sequential native forwards for horizon-specific families. Inputs are synthetic standardized tensors; no test loader or labels are accessed.}
\parbox{\textwidth}{\footnotesize Train (GPU h) is the sum of native per-epoch training times recorded in the frozen logs and excludes validation, testing and unlogged orchestration overhead.}
\end{table}

\end{document}
